% 编译：在仓库根目录执行以下命令两遍。
% mkdir -p output/pdf
% xelatex -halt-on-error -output-directory=output/pdf assignments/assignment11/assignment11-sol.tex
% 中文字体使用 TeX Live / MacTeX 自带的 Fandol 字体。
\documentclass[11pt,a4paper,fontset=fandol]{ctexart}
\usepackage[margin=2.1cm]{geometry}
\usepackage{amsmath,amssymb,mathtools}
\usepackage{fancyhdr,needspace}
\usepackage[hidelinks]{hyperref}
\hypersetup{pdftitle={线性代数第十一次作业},pdfsubject={MATH1205：线性方程组、基础解系与矩阵的秩}}
\setlength{\parindent}{0pt}
\setlength{\parskip}{0.25em}
\setlength{\headheight}{15pt}
\renewcommand{\arraystretch}{1.1}
\linespread{1.08}
\pagestyle{fancy}
\fancyhf{}
\fancyhead[L]{MATH1205\quad 线性代数}
\fancyhead[R]{第十一次作业}
\fancyfoot[C]{\thepage}
\renewcommand{\headrulewidth}{0.3pt}
\ctexset{section={format=\large\bfseries,beforeskip=1.6ex,afterskip=0.7ex}}
\newcommand{\problem}[1]{\Needspace{6\baselineskip}\section*{#1}}
\newcommand{\ans}[1]{\boxed{#1}}
\newcommand{\rank}{\operatorname{r}}
\newcommand{\Null}{\operatorname{Ker}}
\newcommand{\Col}{\operatorname{Col}}
\newcommand{\R}{\mathbb R}

\begin{document}
\setlength{\abovedisplayskip}{6pt plus 2pt minus 2pt}
\setlength{\belowdisplayskip}{6pt plus 2pt minus 2pt}
\setlength{\abovedisplayshortskip}{4pt plus 1pt minus 1pt}
\setlength{\belowdisplayshortskip}{4pt plus 1pt minus 1pt}
\begin{center}
{\LARGE\bfseries 线性代数第十一次作业}\par
\vspace{0.25em}
{\small MATH1205}
\end{center}
\vspace{-0.3em}
按原题号作答，覆盖所给 PDF 中的全部小问。以下均在实数域上讨论。

\problem{第 24 题：求方程组的全部解}
\textbf{(1)} 系数矩阵经初等行变换化为
\[
\begin{aligned}
A&=\begin{pmatrix}2&-1&3\\1&3&2\\3&-5&4\\1&17&4\end{pmatrix}
\xrightarrow{R_1\leftrightarrow R_2}
\begin{pmatrix}1&3&2\\2&-1&3\\3&-5&4\\1&17&4\end{pmatrix}\\[2pt]
&\xrightarrow[\ R_3\leftarrow R_3-3R_1,\ R_4\leftarrow R_4-R_1\ ]{R_2\leftarrow R_2-2R_1}
\begin{pmatrix}1&3&2\\0&-7&-1\\0&-14&-2\\0&14&2\end{pmatrix}
\sim\begin{pmatrix}1&3&2\\0&7&1\\0&0&0\\0&0&0\end{pmatrix}.
\end{aligned}
\]
故 $\rank(A)=2$，方程组等价于 $x+3y+2z=0,\ 7y+z=0$。
令 $y=t$，得 $x=11t,\ z=-7t$。因此一个基础解系及全部解分别为
\[
\ans{\eta=\begin{pmatrix}11\\1\\-7\end{pmatrix},\qquad
\begin{pmatrix}x\\y\\z\end{pmatrix}=t\eta,\quad t\in\R.}
\]

\textbf{(3)} 前两个方程的“第一式的两倍减第二式”给出
\[
x+4y-3z=2\cdot1-4=-2,
\]
而第三个方程要求 $x+4y-3z=7$，二者矛盾。因此
\[
\ans{\text{该方程组无解，解集为 }\varnothing.}
\]

\problem{第 26 题：基础解系与公共非零解}
\textbf{(1)} 方程组 (I) 为 $x_1+x_2=0,\ x_2-x_4=0$。
令 $x_2=s,\ x_3=t$，其通解为
\[
x=s\begin{pmatrix}-1\\1\\0\\1\end{pmatrix}
+t\begin{pmatrix}0\\0\\1\\0\end{pmatrix},\qquad s,t\in\R.
\]
两个系数向量线性无关，故它们构成 (I) 的一个基础解系。

\textbf{(2)} 方程组 (II) 的任一解为
\[
x=c_1\begin{pmatrix}0\\1\\1\\0\end{pmatrix}
+c_2\begin{pmatrix}-1\\2\\2\\1\end{pmatrix}
=\begin{pmatrix}-c_2\\c_1+2c_2\\c_1+2c_2\\c_2\end{pmatrix}.
\]
代入 (I)，两式都化为 $c_1+c_2=0$。故存在公共非零解，且全部为
\[
\ans{x=t\begin{pmatrix}-1\\1\\1\\1\end{pmatrix},\qquad t\in\R\setminus\{0\}.}
\]

\newpage
\problem{第 28 题 (1)：含参数方程组的解的分类}
方程组为
\[
\begin{cases}
ax+y+z=4,\\
x+by+z=3,\\
x+2by+z=4.
\end{cases}
\]
第三式减第二式得 $by=1$，故 $b=0$ 时无解。
当 $b\ne0$ 时，原方程组等价于
\[
y=\frac1b,\qquad x+z=2,\qquad (a-1)x=2-\frac1b.
\]
于是有以下三种情形。

\textbf{唯一解：} $a\ne1,\ b\ne0$ 时，
\[
\ans{(x,y,z)=\left(\frac{2b-1}{b(a-1)},\ \frac1b,\
2-\frac{2b-1}{b(a-1)}\right).}
\]

\textbf{无穷多解：} $a=1,\ b=\tfrac12$ 时，$y=2,\ x+z=2$，全部解为
\[
\ans{\begin{pmatrix}x\\y\\z\end{pmatrix}
=\begin{pmatrix}0\\2\\2\end{pmatrix}
+t\begin{pmatrix}1\\0\\-1\end{pmatrix},\qquad t\in\R.}
\]

\textbf{无解：} $b=0$，或 $a=1,\ b\notin\{0,\tfrac12\}$ 时无解。
后者因为 $0=2-1/b$ 不成立。三类条件互斥且覆盖所有 $a,b$。

\problem{第 29 题：范德蒙德型方程组}
\textbf{(1)} 若 $a_1,a_2,a_3,a_4$ 两两不同，则增广矩阵的行列式为
\[
\det\begin{pmatrix}
1&a_1&a_1^2&a_1^3\\
1&a_2&a_2^2&a_2^3\\
1&a_3&a_3^2&a_3^3\\
1&a_4&a_4^2&a_4^3
\end{pmatrix}
=\prod_{1\le i<j\le4}(a_j-a_i)\ne0.
\]
因此增广矩阵的秩为 $4$，而系数矩阵只有三列，秩至多为 $3$。
两者秩不等，故原方程组\textbf{无解}。

\textbf{(2)} 当 $a_1=a_3=a,\ a_2=a_4=-a,\ a\ne0$ 时，只有两个不同的方程：
\[
\begin{cases}x_1+ax_2+a^2x_3=a^3,\\x_1-ax_2+a^2x_3=-a^3.\end{cases}
\quad\Longleftrightarrow\quad
\begin{cases}x_1+a^2x_3=0,\\x_2=a^2.\end{cases}
\]
所以，暂不使用给定的两个解时，通解为
\[
x=\begin{pmatrix}0\\a^2\\0\end{pmatrix}
+s\begin{pmatrix}-a^2\\0\\1\end{pmatrix},\qquad s\in\R.
\]
题设另给 $\eta_1=c_1(-1,1,1)^T,\ \eta_2=c_1(1,1,-1)^T$ 都是解。
由 $\eta_1$ 的第二分量得 $c_1=a^2\ne0$，再代入第一式得
$c_1(a^2-1)=0$，故 $a^2=1,\ c_1=1$。
此时 $\eta_2$ 也满足方程组。因而在完整题设下，全部解为
\[
\ans{x=\begin{pmatrix}0\\1\\0\end{pmatrix}
+s\begin{pmatrix}-1\\0\\1\end{pmatrix},\qquad s\in\R.}
\]
等价地，可利用给定的两个解写成
$x=\eta_1+t(\eta_2-\eta_1)=(-1+2t,1,1-2t)^T$，$t\in\R$。

\newpage
\problem{第 30 题：判断命题并说明理由}
\textbf{(1) 错误。} 与基础解系等价的向量组只保证张成同一解空间，
不保证线性无关。例如
\[
\eta_1,\ \eta_2,\ \eta_3,\ \eta_1+\eta_2
\]
与 $\eta_1,\eta_2,\eta_3$ 等价，却线性相关，因而不是基础解系。
若另外限定等价组也恰有三个向量，则该组才必为基础解系。

\medskip
\textbf{(2) 错误。} $m<n$ 只能保证\emph{有解时}有无穷多解，不能保证相容。
例如取
\[
A=\begin{pmatrix}1&0&0\\0&0&0\end{pmatrix},\qquad
\beta=\begin{pmatrix}0\\1\end{pmatrix}.
\]
此时 $m=2<n=3,\ \beta\ne0$，但方程组包含 $0=1$，无解。

\medskip
\textbf{(3) 错误。} $\rank(A)=n$ 只能保证\emph{有解时}唯一。
例如取
\[
A=\begin{pmatrix}1\\0\end{pmatrix},\qquad
\beta=\begin{pmatrix}0\\1\end{pmatrix}.
\]
此时 $n=1,\ \rank(A)=1=n$，但 $Ax=\beta$ 要求 $x=0$ 且 $0=1$，故无解。

\medskip
\textbf{(4) 正确。} 因为增广矩阵仍只有 $m$ 行，所以
\[
m=\rank(A)\le\rank(A,\beta)\le m.
\]
故 $\rank(A,\beta)=\rank(A)$，对任意 $\beta\in\R^m$ 都有解。

\medskip
\textbf{(5) 正确。} 先证明 $\rank(A^TA)=\rank(A)$。
显然 $Ax=0$ 推出 $A^TAx=0$；反之，若 $A^TAx=0$，则
\[
0=x^TA^TAx=(Ax)^T(Ax)=\sum_{i=1}^{m}(Ax)_i^2,
\]
由实数平方和为零可得 $Ax=0$。两个齐次方程组同解，且未知数个数相同，
故由解空间维数公式得 $\rank(A^TA)=\rank(A)$。于是
\[
\begin{aligned}
\rank(A^TA)
&\le\rank(A^TA,A^T\beta)\\
&=\rank\bigl(A^T(A,\beta)\bigr)
\le\rank(A^T)=\rank(A)=\rank(A^TA).
\end{aligned}
\]
所以系数矩阵与增广矩阵同秩，$(A^TA)x=A^T\beta$ 必有解。

\medskip
\textbf{(6) 错误。} $Ax=0$ 只有零解，等价于 $\rank(A)=n$，
仍不能保证非齐次方程组有解。第 (3) 问的反例同样适用：
$Ax=0$ 只有零解，而 $Ax=\beta$ 无解。

\medskip
因此，六个命题的判断依次为
\[
\ans{\text{错，错，错，对，对，错。}}
\]

\newpage
\problem{第 32 题：用代数余子式构造基础解系}
已知 $\det A=0$，且第 $i$ 行某元素的代数余子式 $A_{ij}\ne0$。
这说明存在一个非零的 $n-1$ 阶子式，故
\[
n-1\le\rank(A)<n,\qquad\text{即}\quad\rank(A)=n-1.
\]
因此 $Ax=0$ 的解空间维数为 $n-\rank(A)=1$。

令 $v=(A_{i1},A_{i2},\ldots,A_{in})^T$。
注意 $v$ 是伴随矩阵 $\operatorname{adj}(A)$ 的第 $i$ 列。
由伴随矩阵恒等式
\[
A\operatorname{adj}(A)=(\det A)I_n=0
\]
可知 $Av=0$。又因 $v$ 的第 $j$ 个分量 $A_{ij}\ne0$，故 $v\ne0$。
一维解空间中的任意一个非零向量都构成一组基，所以
\[
\ans{\bigl\{(A_{i1},A_{i2},\ldots,A_{in})^T\bigr\}
\text{ 是 }Ax=0\text{ 的一个基础解系。}}
\]

\problem{第 34 题：非齐次解与基础解系的关系}
已知 $A\xi=\beta\ne0$，且 $\eta_1,\ldots,\eta_r$ 是 $Ax=0$ 的一个基础解系。

\textbf{(1)} 设
\[
k_1\eta_1+\cdots+k_r\eta_r+k\xi=0.
\]
左乘 $A$ 得 $k\beta=0$。由 $\beta\ne0$，得 $k=0$。
再由 $\eta_1,\ldots,\eta_r$ 线性无关，得 $k_1=\cdots=k_r=0$。
故 $\eta_1,\ldots,\eta_r,\xi$ \textbf{线性无关}。

\textbf{(2)} 设 $c_0\xi+\sum_{i=1}^r c_i(\eta_i+\xi)=0$，整理为
\[
\sum_{i=1}^r c_i\eta_i+\left(c_0+\sum_{i=1}^r c_i\right)\xi=0.
\]
由第 (1) 问，$c_i=0\ (1\le i\le r)$，且 $c_0+\sum_i c_i=0$，
所以 $c_0=0$。故 $\xi,\eta_1+\xi,\ldots,\eta_r+\xi$ \textbf{线性无关}。

\textbf{(3)} 对任一解 $\gamma$，有 $A(\gamma-\xi)=0$。
由于 $\eta_1,\ldots,\eta_r$ 是基础解系，存在 $t_1,\ldots,t_r$ 使
\[
\gamma=\xi+\sum_{i=1}^r t_i\eta_i.
\]
令 $c_i=t_i\ (1\le i\le r)$，$c_0=1-\sum_{i=1}^r t_i$，便得到
\[
\ans{\gamma=c_0\xi+\sum_{i=1}^r c_i(\eta_i+\xi),\qquad
c_0+c_1+\cdots+c_r=1.}
\]
反之，任何满足上述系数条件的组合都满足
$A\gamma=(c_0+\cdots+c_r)\beta=\beta$，故此表示恰好给出全部解。

\newpage
\problem{第 43 题：实矩阵及其 Gram 矩阵的秩}
设 $A\in\R^{m\times n}$。

\textbf{(1)} 若 $Ax=0$，则显然 $A^TAx=0$。
反之，若 $A^TAx=0$，则
\[
0=x^TA^TAx=(Ax)^T(Ax)=\lVert Ax\rVert^2.
\]
因 $Ax$ 为实向量，范数平方为零当且仅当 $Ax=0$。故
\[
\ans{\Null(A)=\Null(A^TA),\quad
Ax=0\text{ 与 }A^TAx=0\text{ 同解。}}
\]

\textbf{(2)} 第 (1) 问中的两个方程组都有 $n$ 个未知数，由秩与零空间维数公式，
\[
\rank(A)=n-\dim\Null(A)
=n-\dim\Null(A^TA)=\rank(A^TA).
\]
将同一结论用于实矩阵 $A^T\in\R^{n\times m}$，得
\[
\rank(A^T)=\rank\bigl((A^T)^TA^T\bigr)=\rank(AA^T).
\]
再结合矩阵的行秩等于列秩，即 $\rank(A)=\rank(A^T)$，得到
\[
\ans{\rank(A)=\rank(A^TA)=\rank(A^T)=\rank(AA^T).}
\]

\problem{第 49 题：矩阵幂的秩最终稳定}
设 $A$ 为 $n$ 阶矩阵，约定 $A^0=I_n$，记
\[
V_k=\Col(A^k)\qquad(k=0,1,2,\ldots).
\]
由 $A^{k+1}=A^kA$，$A^{k+1}$ 的每一列都是 $A^k$ 的列的线性组合，故
\[
V_0\supseteq V_1\supseteq V_2\supseteq\cdots,
\qquad \dim V_k=\rank(A^k).
\]
若某个 $k$ 满足 $\dim V_k=\dim V_{k+1}$，则由包含关系有 $V_k=V_{k+1}$。
对两边作用 $A$，得到
\[
V_{k+1}=A(V_k)=A(V_{k+1})=V_{k+2}.
\]
重复此过程可知，从这一步起所有 $V_j$ 相等，秩也不再下降。

现分两种情形：若在 $V_0,V_1,\ldots,V_n$ 中已经出现相邻两项相等，
则必有 $V_n=V_{n+1}$；若此前每一步都严格缩小，则维数至少下降 $n$ 次，故
\[
0\le\dim V_n\le\dim V_0-n=n-n=0.
\]
于是 $V_n=\{0\}$，从而 $V_{n+1}=\{0\}$，仍有 $V_n=V_{n+1}$。
因此总有
\[
\ans{\rank(A^n)=\rank(A^{n+1}).}
\]
\end{document}
